\begin{proof}
\textbf{Theorem (Triangle Inequality).} Let $(V,\langle\cdot,\cdot\rangle)$ be a real inner‑product space and denote the induced norm by $\|v\|=\sqrt{\langle v,v\rangle}$. For all vectors $x,y\in V$ one has
\[
\|x+y\|\le \|x\|+\|y\|.
\]

\medskip
\textbf{1. Degenerate case.}  
If $x=0$ or $y=0$ then $\|x+y\|=\|x\|+\|y\|$ because $\|0\|=0$. Hence the inequality holds with equality.

\medskip
\textbf{2. Squaring the inequality.}  
Assume $x\neq0$ and $y\neq0$. Both sides are non‑negative, so we may square:
\begin{align*}
\|x+y\|^{2}\le (\|x\|+\|y\|)^{2}. \tag{1}
\end{align*}

\medskip
\textbf{3. Expansion of the squared norm.}  
Using bilinearity and symmetry,
\begin{align*}
\|x+y\|^{2}
&=\langle x+y,x+y\rangle\\
&=\langle x,x\rangle+\langle x,y\rangle+\langle y,x\rangle+\langle y,y\rangle\\
&=\|x\|^{2}+2\langle x,y\rangle+\|y\|^{2}. \tag{2}
\end{align*}
(The identity was verified with \texttt{verify\_arithmetic}.)

\medskip
\textbf{4. Cauchy–Schwarz bound for the mixed term.}  
The Cauchy–Schwarz inequality gives
\[
|\langle x,y\rangle|\le \|x\|\,\|y\|.
\]
Hence
\begin{align*}
\langle x,y\rangle\le \|x\|\,\|y\|. \tag{3}
\end{align*}
Multiplying by $2$,
\begin{align*}
2\langle x,y\rangle\le 2\|x\|\,\|y\|. \tag{4}
\end{align*}

\medskip
\textbf{5. Substituting the bound.}  
Insert (4) into (2):
\begin{align*}
\|x+y\|^{2}
&=\|x\|^{2}+2\langle x,y\rangle+\|y\|^{2}\\
&\le \|x\|^{2}+2\|x\|\,\|y\|+\|y\|^{2}\\
&=(\|x\|+\|y\|)^{2}. \tag{5}
\end{align*}
Thus (5) establishes (1).

\medskip
\textbf{6. Taking square roots.}  
Both sides of (1) are non‑negative, so taking square roots preserves the inequality:
\[
\|x+y\|\le \|x\|+\|y\|.
\]

\medskip
\textbf{7. Equality condition.}  
Equality can occur only when equality holds in the Cauchy–Schwarz step (3).  
Cauchy–Schwarz is an equality iff $x$ and $y$ are linearly dependent, i.e. $y=\lambda x$ for some $\lambda\in\mathbb R$.  
From (4) we need $2\langle x,y\rangle=2\|x\|\,\|y\|$, which forces $\lambda\ge0$.  
Thus
\[
\|x+y\|=\|x\|+\|y\|
\iff y=\lambda x\text{ with }\lambda\ge0
\]
(and symmetrically $x=\mu y$ with $\mu\ge0$), including the degenerate case when one vector is zero.

\medskip
\textbf{Conclusion.}  
All cases have been considered and each step verified, so the triangle inequality holds in any real inner‑product space. \(\square\)
\end{proof}